\documentclass[book.tex]{subfiles}

\begin{document}

\chapter{Mathematical Transforms}
\label{chap:transforms}
\noindent\rule{11cm}{0.4pt} \\
\textbf{Prerequisites:}  \autoref{chap:calculus}, \autoref{chap:vectors} \\
\textbf{Difficulty Level:} ** \\
\noindent\rule{11cm}{0.4pt}
\vspace{5mm}

We will at various times have to make use of various transforms that turn one type of equation into another. These transforms will prove especially useful to us as we encounter complex sets of differential equations. Having a toolbox of transforms we can throw at systems will enable us to solve much more sophisticated systems than we could do otherwise.


\section{Fourier Transforms}

The Fourier transform is a mathematical operation of fundamental importance to modern physics and engineering. Let $f$ be an integrable function $f: \mathbb{R} \to \mathbb{C}$. 

\begin{equation}
    \hat{f}(k) = \frac{1}{2\pi} \int_{-\infty}^{\infty} f(x) e^{i x k} dx 
\end{equation}

We will also use the notation

\[ \mathcal{F}(f) = \hat{f} \]

to denote the Fourier transform as an operator $\mathcal{F}$. If you're not familiar with the Fourier transform, this definition may feel quite opaque. Why does it make sense? Let's start by first outlining a few basic properties of the Fourier transform

\begin{align}
    \mathcal{F}(a f(x) + b g(x)) &= \int_{-\infty}^{\infty} (a f(x) + b g(x))e^{-2\pi i x k} dx \\
    &= \frac{a}{2\pi} \int_{-\infty}^{\infty} f(x) e^{i x k} dx + \frac{b}{2\pi} \int_{-\infty}^{\infty} g(x) e^{i x k} dx \\
    &= a \mathcal{F}(f(x)) + b \mathcal{F}(g(x))
\end{align}

Put another way, $\mathcal{F}$ is a linear operator. Fourier transforms have a number of really powerful other properties. Let's consider another such property. Let $f$ be a function and $f'$ its derivative. Then

\begin{align}
    \mathcal{F}(f')(k) &= \frac{1}{2\pi} \int_{-\infty}^{\infty} f'(x)e^{ixk} dx
\end{align}

To compute this integral, we can use integration by parts. Let $u = e^{ixk}$ and $dv = f'(x)dx$. Then $du = ik e^{ixk}$ and $v = f(x)$. Then we can write

\begin{align}
    \mathcal{F}(f')(k) &= \frac{1}{2\pi}\int_{-\infty}^{\infty} u dv \\
    &= \frac{1}{2\pi} \left (uv \Big |^{\infty}_{-\infty} - \int _{-\infty}^{\infty} v du \right ) \\
    &= \frac{1}{2\pi} \left (e^{ixk} f(x) \Big |^{\infty} _{-\infty} - \int _{-\infty}^{\infty} f(x) ike^{ixk} dx \right ) \\
\end{align}

We assume that $f(x)$ goes to 0 at both infinities, so we are left with

\begin{align}
\mathcal{F}(f')(k) &= \frac{-ik}{2\pi} \int_{-\infty}^{\infty} f(x) e^{ixk} dx \\
&= -ik \mathcal{F}(f)(k)
\end{align}

That is, the Fourier transform converts a derivative into multiplication by a scalar constant.% We will return to this property to solve partial differential equations later in the book.

%\textcolor{red}{TODO: Cover Fourier position/momentum space connection.}

%\textcolor{red}{TODO: Cover Fourier Space on $\mathbb{R}^n$}



\section{Convexity and Legendre Transform}

The Legendre transform is a type of functional transformation that arises in a number of different places in physics. In this section we'll give a brief introduction to the Legendre transform. We'll start by defining a \textit{convex} sets and functions.

Let $U \subseteq \mathbb{R^n}$. We say that $U$ is convex if for all $\vec{v}, \vec{u} \in U$, and for all $t \in [0, 1]$, we have that $t\vec{v} + (1-t)\vec{u} \in U$.

%\textcolor{red}{TODO: Add picture of a convex set.}

A function $f: U \to \mathbb{R}$, where $U \subseteq \mathbb{R}^n$ is said to be convex if for all $\vec{v}, \vec{u} \in U$, and for all $t \in [0, 1]$, we have
\begin{align}
f(t\vec{v} + (1-t)\vec{u}) \leq tf(\vec{v}) + (1-t)f(\vec{u})
\end{align}

%\textcolor{red}{TODO: Add picture of convex function.}

Convex functions arise often in optimization and in physics. Intuitively, a convex function is "well-behaved". The Legendre transform is defined for convex functions as follows.

%\textcolor{red}{TODO: Add a definition of a supremum.}
Let $U \subseteq \mathbb{R}^n$ be a convex set. Let $f: U \to \mathbb{R}$ be a convex function. We defined the Legendre transform of $f$ to be a function $f$* defined by 

\begin{align}
f^*(y) = \mathrm{sup}_{\vec{x} \in U} \left (\vec{y}^T\vec{x} - f(x) \right )
\end{align}

where $\vec{y} \in V \subseteq \mathbb{R}^n$ is a set defined by

\begin{align}
V = \{\vec{y} \in \mathbb{R}^n : f^*(y) < \infty \}
\end{align}

Here $f^*$ is referred to as the convex conjugate of $f$.

The mathematical method of performing this change of variables is called a Legendre Transform. The method is defined according to the following steps:

A function $y=f(x)$ is defined as a list of $x$, $y$, pairs, \textit{i.e.} the plot of $y(x)$ graphs 
\begin{align} 
(x_1,y_1), (x_2,y_2), (x_3,y_3), \dotsc 
\end{align} 
and so on. Equivalently, it is possible to express the same function in terms of the tangent slope $c(x)$ and the corresponding $y$-intercepts, \textit{i.e.} the points 
\begin{align}
(c_1,b_1), (c_2,b_2), (c_3,b_3),\dotsc
\end{align}
For a change $dx$ in the $x$-variable, the change $dy$ is given by:
\begin{align}
dy = \left(\frac{dy}{dx}\right)dx = c(x)dx \notag
\end{align} 
From this point on the curve, we can extend a line back to $x = 0$ to find the $y$-intercept $b(x)$, such that
\begin{align}
y(x) = c(x)x+b(x) \rightarrow b(x) = y(x)-c(x)x \notag
\end{align}
The resulting function for the series of intercepts versus the slopes (\textit{i.e.} $b = b(c)$) contains the same information as $y = y(x)$.

A change in $b(c)$ is given by 
\begin{align}
db &= dy - cdx - xdc = -xdc
\end{align}

As an example, consider the function $y = (x - 2)^2$. We use a Legendre Transform to shift the variable from $x$ to the slope $c$. 

The slope of the function is: 
\begin{align}
c = 2(x-2) \rightarrow x =\frac{c}{2} + 2
\end{align}

The Legendre transform gives us a new function $b(c)$, given by:
\begin{align}
b=y-cx=\frac{c^2}{4} - \frac{c^2}{2}-2c=-\frac{c^2}{4}-2c 
\end{align}

We note that this Legendre transform correctly relates back to the
original function by noting that
\begin{align}
\frac{db}{dc} = -\frac{c}{2} -2 = -x \notag
\end{align}
end{example}
\begin{marginfigure}%
\includegraphics[width=\linewidth]{figures/mathematical_transforms/thermo_legendre.png}
\label{fig:Fig5}
\end{marginfigure}


\section{Laplace Transform}

The Laplace transform provides a technique for transforming a real-valued function into a complex-valued function in a structured way. Like the Fourier transform, it will enable us to simplify certain complex mathematical operations.

Let's suppose that we have a function $f: \mathbb{R}^+ \to \mathbb{R}$. We then define

\begin{align}
    \mathcal{L}\{f\}(s) &= \int_0^{\infty} f(t) e^{-s t} dt 
\end{align}
where $s$ is a complex number. That means $L\{f\}$ has type $\mathbb{C} \to \mathbb{C}$. Structurally, you might note the similarities between the Fourier transform and the Laplace transform. And indeed, the two share similar properties. To start, $\mathcal{L}$ is a linear operator

\begin{align}
    \mathcal{L}\{af + bg\} &= \int_{0}^{\infty} (a f(t) + b g(t)) e^{-st} dt \\
    &= a\int_{0}^{\infty} f(t)e^{-st} dt + b \int_{0}^{\infty} g(t) e^{-st} dt \\
    &= a\mathcal{L}\{f\} + b \mathcal{L}\{g\}
\end{align}

%TODO: Add some more exercises here.

\subsection{Exercises}
\begin{enumerate}
    \item Let $f(x) = e^x$.
    \begin{enumerate}
        \item Prove that $f$ is a convex function. 
        \item Show that the convex conjugate $f^*(y) = y \ln y - y$. 
        \item Show that the domain of $f^*$ is the set $\{y \geq 0\ |\ y \in \mathbb{R}\}$.
    \end{enumerate}
\end{enumerate}

\end{document}